% --- Archon physics formalization source begin ---
% archon:physics
% archon:covers IPhO2026Problems/problem_IPhO_2026_1_B_1.lean
% archon:source-report reports/ipho_2026/problem_IPhO_2026_1_B_1.source.json
% archon:problem-id IPhO_2026_1
% archon:part-id B.1

\chapter{Physics problem IPhO\_2026\_1\_B\_1}
\label{ch:IPhO2026Problems_problem_IPhO_2026_1_B_1}

\paragraph{Problem source.}
At one instant a positron and an electron, each of mass m and charges of equal
magnitude and opposite sign, are separated by 100*a\_0.  Their velocities are
antiparallel and perpendicular to their separation.  Each particle has angular
momentum of magnitude mu*hbar about the center of mass.  The system is isolated,
classical, non-relativistic, and has only electrostatic interaction.  The Bohr
radius is a\_0 = 4*pi*epsilon\_0*hbar\textasciicircum{}2/(m*e\textasciicircum{}2), and k = 1/(4*pi*epsilon\_0).

Current subquestion:
For mu = 4 the pair is bound. Find the maximum electron-positron separation in units of a\_0.

\paragraph{Current subquestion.}
For mu = 4 the pair is bound. Find the maximum electron-positron separation in units of a\_0.

\paragraph{Recorded answer/context.}
r\_max = (1600/9)*a\_0.

\paragraph{Figure/image path.}
/root/proposal\_for\_physic/science-mango/ipho\_2026\_source/image/T1\_page-2.png

\paragraph{Formalization target.}
create a compiling Lean file with sorry bodies at `IPhO2026Problems/problem\_IPhO\_2026\_1\_B\_1.lean`. The Lean declarations must preserve the physical quantities, dimensions or dimensional roles, figure labels, governing-law hypotheses, and final relation expressed by this problem.
Use Mathlib/Physlib names found through LeanExplore where available. If a physics API is missing, introduce faithful local abstractions rather than scalar placeholder aliases.

\begin{theorem}[Physics formalization target]
  \label{thm:physics:IPhO_2026_1_B_1:target}
  \lean{IPhO2026Problem1B1.maximum_separation_for_mu_four}
  \uses{def:physics:IPhO_2026_1_B_1:aux001, def:physics:IPhO_2026_1_B_1:aux002, def:physics:IPhO_2026_1_B_1:aux003, def:physics:IPhO_2026_1_B_1:aux004, def:physics:IPhO_2026_1_B_1:aux005, def:physics:IPhO_2026_1_B_1:aux006, def:physics:IPhO_2026_1_B_1:aux007, def:physics:IPhO_2026_1_B_1:aux008, def:physics:IPhO_2026_1_B_1:aux009, def:physics:IPhO_2026_1_B_1:aux010, def:physics:IPhO_2026_1_B_1:aux011, def:physics:IPhO_2026_1_B_1:aux012, def:physics:IPhO_2026_1_B_1:aux013}
  For μ = 4, the maximum electron--positron separation is (1600 / 9) a₀.
\end{theorem}
\begin{proof}
  Use the typed geometry, governing-law, branch, and measurement interfaces listed in the declaration index.  Eliminate the intermediate readouts in their physical order and then evaluate the requested exact or uncertainty relation.
\end{proof}
\section*{Formal declaration index}
The following blocks record the typed carriers and bridge declarations used by the target.  Their dependency lists follow the declaration-level model in the covered Lean file.

\begin{definition}[Declaration LengthQuantity]
  \label{def:physics:IPhO_2026_1_B_1:aux001}
  \lean{IPhO2026Problem1B1.LengthQuantity}
  A scalar length measured in an arbitrary but fixed system of units.
\end{definition}
\begin{proof}
  This is the typed definition of the stated physical carrier; its fields or defining equation provide the claimed data.
\end{proof}

\begin{definition}[Declaration MassQuantity]
  \label{def:physics:IPhO_2026_1_B_1:aux002}
  \lean{IPhO2026Problem1B1.MassQuantity}
  A scalar mass measured in an arbitrary but fixed system of units.
\end{definition}
\begin{proof}
  This is the typed definition of the stated physical carrier; its fields or defining equation provide the claimed data.
\end{proof}

\begin{definition}[Declaration ChargeQuantity]
  \label{def:physics:IPhO_2026_1_B_1:aux003}
  \lean{IPhO2026Problem1B1.ChargeQuantity}
  A signed electric charge measured in an arbitrary but fixed system of units.
\end{definition}
\begin{proof}
  This is the typed definition of the stated physical carrier; its fields or defining equation provide the claimed data.
\end{proof}

\begin{definition}[Declaration PositionQuantity]
  \label{def:physics:IPhO_2026_1_B_1:aux004}
  \lean{IPhO2026Problem1B1.PositionQuantity}
  A three-dimensional position vector.
\end{definition}
\begin{proof}
  This is the typed definition of the stated physical carrier; its fields or defining equation provide the claimed data.
\end{proof}

\begin{definition}[Declaration VelocityQuantity]
  \label{def:physics:IPhO_2026_1_B_1:aux005}
  \lean{IPhO2026Problem1B1.VelocityQuantity}
  A three-dimensional velocity vector.
\end{definition}
\begin{proof}
  This is the typed definition of the stated physical carrier; its fields or defining equation provide the claimed data.
\end{proof}

\begin{definition}[Declaration ActionQuantity]
  \label{def:physics:IPhO_2026_1_B_1:aux006}
  \lean{IPhO2026Problem1B1.ActionQuantity}
  Action, and hence angular momentum, with dimension M L² T⁻¹.
\end{definition}
\begin{proof}
  This is the typed definition of the stated physical carrier; its fields or defining equation provide the claimed data.
\end{proof}

\begin{definition}[Declaration EnergyQuantity]
  \label{def:physics:IPhO_2026_1_B_1:aux007}
  \lean{IPhO2026Problem1B1.EnergyQuantity}
  Energy with dimension M L² T⁻².
\end{definition}
\begin{proof}
  This is the typed definition of the stated physical carrier; its fields or defining equation provide the claimed data.
\end{proof}

\begin{definition}[Declaration CoulombConstantQuantity]
  \label{def:physics:IPhO_2026_1_B_1:aux008}
  \lean{IPhO2026Problem1B1.CoulombConstantQuantity}
  Coulomb's constant, with dimension M L³ T⁻² C⁻².
\end{definition}
\begin{proof}
  This is the typed definition of the stated physical carrier; its fields or defining equation provide the claimed data.
\end{proof}

\begin{definition}[Declaration VacuumPermittivityQuantity]
  \label{def:physics:IPhO_2026_1_B_1:aux009}
  \lean{IPhO2026Problem1B1.VacuumPermittivityQuantity}
  Vacuum permittivity, with dimension M⁻¹ L⁻³ T² C².
\end{definition}
\begin{proof}
  This is the typed definition of the stated physical carrier; its fields or defining equation provide the claimed data.
\end{proof}

\begin{definition}[Declaration ParticleState]
  \label{def:physics:IPhO_2026_1_B_1:aux010}
  \lean{IPhO2026Problem1B1.ParticleState}
  \uses{def:physics:IPhO_2026_1_B_1:aux002, def:physics:IPhO_2026_1_B_1:aux003, def:physics:IPhO_2026_1_B_1:aux004, def:physics:IPhO_2026_1_B_1:aux005, def:physics:IPhO_2026_1_B_1:aux006}
  The state variables displayed for one of the two particles in Fig. 1b.
\end{definition}
\begin{proof}
  This is the typed definition of the stated physical carrier; its fields or defining equation provide the claimed data.
\end{proof}

\begin{definition}[Declaration CoulombPairSystem]
  \label{def:physics:IPhO_2026_1_B_1:aux011}
  \lean{IPhO2026Problem1B1.CoulombPairSystem}
  \uses{def:physics:IPhO_2026_1_B_1:aux001, def:physics:IPhO_2026_1_B_1:aux003, def:physics:IPhO_2026_1_B_1:aux006, def:physics:IPhO_2026_1_B_1:aux007, def:physics:IPhO_2026_1_B_1:aux008, def:physics:IPhO_2026_1_B_1:aux009, def:physics:IPhO_2026_1_B_1:aux010}
  All named physical quantities used in the electron--positron model. The scalar fields are readouts in one fixed unit system, while their types retain their physical dimensions. mu is the dimensionless angular-momentum factor.
\end{definition}
\begin{proof}
  This is the typed definition of the stated physical carrier; its fields or defining equation provide the claimed data.
\end{proof}

\begin{definition}[Declaration turningPointEnergyReadout]
  \label{def:physics:IPhO_2026_1_B_1:aux012}
  \lean{IPhO2026Problem1B1.turningPointEnergyReadout}
  \uses{def:physics:IPhO_2026_1_B_1:aux001, def:physics:IPhO_2026_1_B_1:aux011}
  The classical effective energy of the pair at a radial turning point of separation r. For either particle, ℓ = m (r/2) v, so the two kinetic energies sum to 4 ℓ² / (m r²). The electrostatic potential energy of the opposite charges is -k e² / r.
\end{definition}
\begin{proof}
  This is the typed definition of the stated physical carrier; its fields or defining equation provide the claimed data.
\end{proof}

\begin{definition}[Declaration CoulombPairLaws]
  \label{def:physics:IPhO_2026_1_B_1:aux013}
  \lean{IPhO2026Problem1B1.CoulombPairLaws}
  \uses{def:physics:IPhO_2026_1_B_1:aux011, def:physics:IPhO_2026_1_B_1:aux012}
  The figure readouts and physical laws for the isolated classical Coulomb pair. The two energy equations are the usable mathematical consequences of the classical, non-relativistic, electrostatic-only, isolated-system assumptions. The final numerical value of maximumSeparation is deliberately not a field.
\end{definition}
\begin{proof}
  This is the typed definition of the stated physical carrier; its fields or defining equation provide the claimed data.
\end{proof}
% --- Archon physics formalization source end ---
